\section{Introduction}
\label{sec:intro}

Text-to-image attacks can produce an unsafe image while changing the visual
content that was requested. A nude sculpture, a partly exposed illustration,
and an explicit drawn figure in a specified pose can all satisfy a broad
unsafe-content criterion. They demonstrate different capabilities. For an
attacker seeking a particular depiction, success requires the intended figure
material, anatomical detail, and body configuration to occur together. This
paper studies that joint requirement and introduces a contextual embedding
attack that produces modern drawn figures with high exposure in multiple
requested poses.

Prior work explores token substitution, artistic reframing, material
substitution, and semantic contexts to alter how a text-to-image system
interprets a request \citep{sneakyprompt,aestheticcamouflage}. A reformulation
can also alter the body that appears in the output. Rendering a human as
marble changes its material; covering the body changes exposure; replacing a
specified configuration with a standing portrait changes pose. These are
substantive differences between the requested content and the generated
result. Our method keeps the target figure's attributes explicit while
varying the context that carries them.

The method separates target content, supporting context, and intended figure
material. Contextual embedding specifies the relation between the target
content and the task represented by the surrounding scene. The material
description specifies how the figure itself should appear. This separation
allows a photographic outer scene to contain a drawn human, as in a
photographed manga page. Production documents, displayed artworks, and
printed pages provide different contexts for this relation. The generation
objective is to retain the requested figure attributes within the resulting
composition. Our outputs demonstrate the joint attainment of a modern drawn
material, maximum exposure under our criteria, and sexualized posture across
different body configurations.

Material and context ablations examine the conditions under which the target
survives. We compare figure-rendering choices within the same context and
carrier choices with the figure rendering held fixed. Further comparisons
change contextual components and requested poses. These experiments connect
the construction to its outputs: a component may affect image return, the
attributes retained in the image, or both. Separating the figure from its
carrier is essential to interpreting these effects. A changed frame or paper
surface can accompany a preserved drawing style, while an apparently
successful generation can lose the requested exposure or configuration.

Evaluating this capability requires more detail than an aggregate attack
success rate (ASR). The rate measures success under a chosen criterion;
it does not describe the mix of materials, exposure levels, and poses that
contribute to that success. We therefore assess these three attributes
independently and test their joint agreement with the target. Material
assessment concerns the depicted body, exposure assessment concerns visible
anatomical features, and pose assessment concerns body configuration and
interaction. The resulting profile explains what each output achieves.
Reporting joint attainment alongside its components makes attack success
interpretable under explicit visual requirements.

\paragraph{Contributions.}
\textbf{(i) A contextual embedding attack.} A prompt-only black-box method
that relates target content to a supporting context while preserving explicit
requirements for the figure's material, exposure, and pose.
\textbf{(ii) Material and context ablations.} Comparisons of figure rendering,
external carriers, and contextual components, supported by outputs that
jointly attain modern drawn material, maximum exposure, and sexualized posture.
\textbf{(iii) A decoupled evaluation framework.} Independent assessment of
material, exposure, and pose, combined with target matching to identify the
specific capability behind an attack's success rate.
